% eksperymentalne tłumaczenie na włoski

%

\pol{\section{Twierdzenia Ptolemeusza, Carnota, Caseya}}
\ita{\section{Teoremi di Tolomeo, Carnot e Casey}}

\pol{Klaudiusz Ptolemeusz będzie astronomem, matematykiem i~geografem pochodzenia greckiego.}
\ita{Claudio Tolomeo fu un astronomo, matematico e geografo di origine greca.}
\index[persons]{Ptolemeusz, Klaudiusz}%
\pol{Urodzon w Tebaidzie (około roku 100), wykształci się i będzie działać w~Aleksandrii; tam też umrze około roku 170.}
\ita{Nato nella Tebaide intorno all'anno 100, si formò e operò ad Alessandria, dove morì intorno al 170.}
\pol{Napisze po grecku Μαθηματικὴ Σύνταξις, traktat w trzynastu księgach znany lepiej jako \emph{Almagest} zawierający kompendium wiedzy astronomicznej oraz matematyczny wykład teorii geocentrycznej, ale również:}
\ita{Scrisse in greco la Μαθηματικὴ Σύνταξις, un trattato in tredici libri meglio noto come \emph{Almagesto}, contenente un compendio delle conoscenze astronomiche e una trattazione matematica della teoria geocentrica, ma anche:}

\begin{theorem}[\pol{Ptolemeusza, 140 r.n.e.}\ita{di Tolomeo, 140 d.C.}]
\label{twierdzenie_ptolemeusza}%
\pol{\index{nierówność!Ptolemeusza}}%
\ita{\index{disuguaglianza!di Tolomeo}}%
\pol{\index{twierdzenie!Ptolemeusza}}%
\ita{\index{teorema!di Tolomeo}}%
    \pol{W czworokącie wypukłym $ABCD$ zachodzi}
    \ita{Per ogni quadrilatero convesso $ABCD$ vale}
    \begin{equation}
        |AC| \cdot |BD| \le |AB| \cdot |CD| + |BC| \cdot |AD|,
    \end{equation}
    \pol{z równością wtedy i tylko wtedy, gdy na czworokącie $ABCD$ można opisać okrąg.}
    \ita{con uguaglianza se e solo se il quadrilatero $ABCD$ è ciclico.}
    \pol{Założenie o wypukłości można pominąć, czwórka punktów nie musi nawet leżeć w~jednej płaszczyźnie -- ale wtedy równość zachodzi też wtedy, kiedy punkty są współliniowe.}
    \ita{L'ipotesi di convessità può essere omessa; i quattro punti non devono nemmeno appartenere allo stesso piano, ma in tal caso l'uguaglianza vale anche quando sono allineati.}
\end{theorem}

\pol{Ptolemeusz udowodni równość, a nie nierówność, ale nazwa się przyjmie.}
\ita{Tolomeo dimostrò l'uguaglianza e non la disuguaglianza, ma il nome rimase comunque in uso.}
\pol{To wystarczy mu, żeby opracować ,,tablice cięciw'' (równoważne tablicom wartości funkcji trygonometrycznych), potrzebne do celów astronomicznych.}
\ita{Questo gli bastò per costruire le \emph{tavole delle corde} (equivalenti alle tavole trigonometriche), necessarie per le applicazioni astronomiche.}
\pol{\index{tablice cięciw}}%
\ita{\index{tavole delle corde}}%
\pol{Wcześniejsze tablice Hipparchosa z~Nikei opisywały tylko wielokrotności kąta miary $\pi/24$.}
\ita{Le precedenti tavole di Ipparco di Nicea descrivevano soltanto multipli dell'angolo $\pi/24$.}
\index[persons]{Hipparchos z Nikei}%
% TODO: {Thurston, Hugh (1996), Early Astronomy, Springer, ISBN 978-0-387-94822-5, strony 235-236}

\pol{O twierdzeniu Ptolemeusza piszą Guzicki \cite[s. 378-379, 407-419]{guzicki_2021}, Bogdańska, Neugebauer \cite[s. 62, 63]{neugebauer_2018}, Audin \cite[s. 108]{audin_2003}.}
\ita{Del teorema di Tolomeo scrivono Guzicki \cite[p. 378-379, 407-419]{guzicki_2021}, Bogdańska e Neugebauer \cite[p. 62, 63]{neugebauer_2018}, Audin \cite[p. 108]{audin_2003}.}
\pol{Dowód inwersyjny przedstawią Eves \cite[s. 132]{eves1_1972} i Waldemar Pompe w~$\Delta_{95}^1$.}
\ita{Una dimostrazione tramite inversione sarà presentata da Eves \cite[p. 132]{eves1_1972} e da Waldemar Pompe in $\Delta_{95}^1$.}
\index[persons]{Pompe, Waldemar}%
% https://www.deltami.edu.pl/1995/01/jeszcze-raz-o-nierownosci-ptolemeusza/
\pol{Krzysztof Apt w $\Delta_{24}^8$ zaprezentuje ,,dowód bez słów'' autorstwa Williama Derricka i~Jamesa Hersteina \cite{derrick_2012}.}
\ita{Krzysztof Apt presenterà in $\Delta_{24}^8$ una \emph{dimostrazione senza parole} di William Derrick e James Herstein \cite{derrick_2012}.}
% https://www.deltami.edu.pl/2024/08/elegancki-dowod-twierdzenia-ptolemeusza/
\index[persons]{Derrick, William}%
\index[persons]{Herstein, James}%
\pol{Jacek Gancarzewicz, Magdalena Staszek w $\Delta_{12}^{12}$ użyją słów, ale nie inwersji.}
\ita{Jacek Gancarzewicz e Magdalena Staszek, in $\Delta_{12}^{12}$, useranno invece le parole, ma non l'inversione.}
% https://www.deltami.edu.pl/2012/12/nierownosc-ptolemeusza/

\pol{Bartłomiej Bzdęga poprosi Czytelnika w $\Delta_{21}^{3}$, by poczuł radość płynącą z zauważenia tego, co się stanie, gdy zastosujemy twierdzenie Ptolemeusza do prostokąta (który jest w oczywisty sposób cykliczny).}
\ita{Bartłomiej Bzdęga inviterà il lettore, in $\Delta_{21}^{3}$, a provare la gioia di osservare ciò che accade applicando il teorema di Tolomeo a un rettangolo (che è ovviamente ciclico).}
% https://www.deltami.edu.pl/2021/03/twierdzenie-ptolemeusza/

\pol{Z twierdzenia Ptolemeusza można wyprowadzić jedno z kilku twierdzeń, które przypisuje się Lazarowi Carnotowi \cite{carnot_1803}:}
\ita{Dal teorema di Tolomeo si può dedurre uno dei vari teoremi attribuiti a Lazare Carnot \cite{carnot_1803}:}
\index[persons]{Carnot, Lazare Nicolas Marguerite}%
% https://ru.wikipedia.org/w/index.php?title=Формула_Карно&oldid=8679639
% В ее доказательстве используется теорема Птолемея.

\begin{theorem}[Carnot, 1803?]
    \pol{\index{twierdzenie!Carnota}}%
    \ita{\index{teorema!di Carnot}}%
    \pol{Niech $ABC$ będzie trójkątem wpisanym w okrąg o środku $O$ i promieniu $R$ oraz opisanym na okręgu o promieniu $r$.}
    \ita{Sia $ABC$ un triangolo inscritto in una circonferenza di centro $O$ e raggio $R$, e circoscritto a una circonferenza di raggio $r$.}
    \pol{Oznaczmy przez $OO_A$ (i analogicznie $OO_B$, $OO_C$) znakowaną odległość punktu $O$ od boku $BC$.}
    \ita{Indichiamo con $OO_A$ (e analogamente con $OO_B$, $OO_C$) la distanza orientata del punto $O$ dal lato $BC$.}
    \pol{Wtedy}
    \ita{Allora}
    \begin{equation}
        OO_A + OO_B + OO_C = R + r.
    \end{equation}
    \pol{(Odległość jest ujemna wtedy i tylko wtedy, gdy cały odcinek leży poza trójkątem).}
    \ita{(La distanza è negativa se e solo se l'intero segmento giace all'esterno del triangolo).}
    \pol{\index{twierdzenie!Carnota}}%
    \ita{\index{teorema!di Carnot}}%
\end{theorem}

\pol{Wynik ten znajduje znowu zastosowanie w dowodzie twierdzenia japońskiego.}
\ita{Questo risultato trova nuovamente applicazione nella dimostrazione del teorema giapponese.}
\index{twierdzenie!japońskie}
\pol{Pisze o nim Zetel \cite[s. 80]{zetel_2020}; Bogdańska i Neugebauer \cite[s. 64]{neugebauer_2018} dla trójkąta ostrokątnego.}
\ita{Ne scrive Zetel \cite[p. 80]{zetel_2020}; Bogdańska e Neugebauer \cite[p. 64]{neugebauer_2018} lo trattano nel caso di un triangolo acutangolo.}

\begin{proposition}[\pol{twierdzenie japońskie}\ita{teorema giapponese}]
    \pol{W wielokącie cyklicznym, który podzielono przekątnymi na trójkąty wpisane w ten sam okrąg, suma promieni okręgów wpisanych w trójkąty jest stała i nie zależy od wybranej triangulacji.}
    \ita{In un poligono ciclico suddiviso mediante diagonali in triangoli inscritti nella stessa circonferenza, la somma dei raggi delle circonferenze inscritte nei triangoli è costante e non dipende dalla triangolazione scelta.}
\end{proposition}

\pol{Zgodnie ze starożytnym zwyczajem japońskich matematyków, twierdzenie to będzie problemem Sangaku wyrytym na tabliczkach zawieszanych w japońskich świątyni ku czci bogów oraz autora około 1800 roku.}
\ita{Secondo l'antica tradizione dei matematici giapponesi, questo teorema apparve come problema \emph{Sangaku}, inciso su tavolette appese nei templi giapponesi in onore delle divinità e dell'autore, intorno al 1800.}
\pol{Piszą o nim Bogdańska, Neugebauer \cite[s. 65]{neugebauer_2018}.}
\ita{Ne scrivono Bogdańska e Neugebauer \cite[p. 65]{neugebauer_2018}.}
\pol{Więcej o sangaku można przeczytać w skrócie pracy Anny Dymek nagrodzonej srebrnym medalem w 33. konkursie uczniowskich prac z~matematyki w 2011 roku (patrz $\Delta_{12}^5$).}
\ita{Maggiori informazioni sui sangaku si possono trovare nel riassunto del lavoro di Anna Dymek, premiato con la medaglia d'argento nel 33º concorso per lavori matematici studenteschi del 2011 (si veda $\Delta_{12}^5$).}

% TODO: https://en.wikipedia.org/wiki/Van_Schooten's_theorem

\begin{theorem}[Caseya, 1866]
\index{twierdzenie!Caseya}%
    \pol{Niech $\Gamma_1$, $\Gamma_2$, $\Gamma_3$, $\Gamma_4$ będą czterema okręgami ponumerowanymi zgodnie z ruchem wskazówek zegara, z których każdy styka się z piątym okręgiem $\Gamma$.}
    \ita{Siano $\Gamma_1$, $\Gamma_2$, $\Gamma_3$, $\Gamma_4$ quattro circonferenze numerate in senso orario, ciascuna tangente a una quinta circonferenza $\Gamma$.}
    \pol{Niech $t_{ij}$ oznacza długość zewnętrznego odcinka stycznego łączącego okręgi $\Gamma_i$, $\Gamma_j$ (jeśli te stykają się z $\Gamma$ obydwa od wewnątrz lub obydwa od zewnątrz) albo długość wewnętrznego odcinka stycznego (w przeciwnym razie).}
    \ita{Sia $t_{ij}$ la lunghezza del segmento tangente esterno che collega le circonferenze $\Gamma_i$ e $\Gamma_j$ (se entrambe toccano $\Gamma$ internamente oppure entrambe esternamente), oppure del segmento tangente interno nel caso contrario.}
    \pol{Wówczas:}
    \ita{Allora:}
    \begin{equation}
        t_{12} \cdot t_{34} + t_{14} \cdot t_{23} = t_{13} \cdot t_{24}.
    \end{equation}
\end{theorem}

% https://en.wikipedia.org/wiki/Casey%27s_theorem
\pol{Twierdzenie poda John Casey (1820-1891), szanowany irlandzki geometra, który razem z Émilem Lemoinem uznawany będzie za współzałożyciela nowoczesnej geometrii trójkątów i okręgów.}
\ita{Il teorema fu enunciato da John Casey (1820--1891), stimato geometra irlandese, considerato insieme a Émile Lemoine uno dei fondatori della moderna geometria dei triangoli e delle circonferenze.}
\index[persons]{Casey, John}%
\index[persons]{Lemoine, Émile}%
\pol{Inny dowód wymyśli Max Zacharias \cite{zacharias_1942}.}
\ita{Una diversa dimostrazione fu trovata da Max Zacharias \cite{zacharias_1942}.}
\index[persons]{Zacharias, Max}%
\pol{Twierdzenie odwrotne do podanego przydaje się w najkrótszym znanym dowodzie twierdzenia Feuerbacha (że okrąg dziewięciu punktów jest styczny do okręgów dopisanych oraz wpisanego).}
\ita{Il teorema inverso è utile nella più breve dimostrazione nota del teorema di Feuerbach (secondo cui la circonferenza dei nove punti è tangente alle circonferenze exinscritte e a quella inscritta).}
\index{okrąg!wpisany}%
\index{okrąg!dopisany}%
\index{twierdzenie!Feuerbacha}%
\index{okrąg!dziewięciu punktów}%

\pol{Znajdziemy je u Bogdańskiej, Neugebauera jako ćwiczenie \cite[s. 105]{neugebauer_2018}.}
\ita{Lo troviamo come esercizio in Bogdańska e Neugebauer \cite[p. 105]{neugebauer_2018}.}
\pol{Wspomni o nim też Joanna Ochremiak w $\Delta_{11}^9$.}
\ita{Lo menziona anche Joanna Ochremiak in $\Delta_{11}^9$.}
% https://www.deltami.edu.pl/2011/09/stowarzyszenie-na-rzecz-edukacji-matematycznej/

%